% !PRINTED 335
\weberpage{335}
%<<CONT
worin $f_k,\Delta_k,J_k$ die Ableitungen nach $x_k$ bedeuten. Für die canonische Form ist dann
\weq{11}{K(y) = r^9 K.}
Um $K(y)$ zu finden, brauchen wir nur die Functionaldeterminante aus den linken Theilen der Formeln (7), (8), (9) zu bilden, woraus sich mit Anwendung einfacher Determinantensätze ergiebt (Bd. I, \S.\,22, \S.\,27):
\[
K(y) = -h^3(h^3+8m^3)^3 \left|\begin{array}{ccc} 1, & y_1^3, & y_1^6 \\ 1, & y_2^3, & y_2^6 \\ 1, & y_3^3, & y_3^6 \end{array}\right|,
\]
und mithin ist
\weq{12}{r^9 K = h^3(h^3+8m^3)^3\,(y_1^3-y_2^3)(y_1^3-y_3^3)(y_2^3-y_3^3).}
Zu bemerken ist noch zu diesen Formeln, dass der in ihnen auftretende Factor $h^3+8m^3$ nicht verschwinden kann, wenn die Curve $f$ keinen singulären Punkt hat. Denn ein singulärer Punkt wird bestimmt durch die drei Gleichungen:
\[
h\,y_1^2 + 2m\,y_2y_3 = 0,\quad h\,y_2^2 + 2m\,y_3y_1 = 0,\quad h\,y_3^2 + 2m\,y_1y_2 = 0,
\]
und diese drei Gleichungen sind mit einander verträglich, wenn $h^3+8m^3=0$. (Ist z.\,B. $h=-2m$, so sind sie befriedigt, wenn $y_1=y_2=y_3$ ist.)

Der Co\"efficient $h$ kann für eine gegebene Form $f$ jeden vorgeschriebenen von Null verschiedenen Werth haben, und erst wenn dieser Werth gegeben ist, ist das Problem bestimmt. Wir wollen der Einfachheit halber jetzt $h=1$ setzen, und erhalten aus (7) bis (9) das Resultat:

Setzt man
\weq{13}{\begin{aligned}
P &= \frac{(2+m^3)m^3f^2+(2-5m^3)mr^2f\Delta+3m^2r^4\Delta^2-r^6J}{(1+8m^3)^2}, \\
Q &= \frac{(1+2m^3)f-6mr^2\Delta}{1+8m^3}, \\
R &= \frac{m^2f+r^2\Delta}{1+8m^3},
\end{aligned}}
so sind $y_1^3,y_2^3,y_3^3$ die Wurzeln der cubischen Gleichung
\weq{14}{u^3 - Qu^2 + Pu - R^3 = 0,}
und die Discriminante dieser cubischen Gleichung ist:
\weq{15}{D_1 = \frac{r^{18}K^2}{(1+8m^3)^6}.}
